% this file is encoded in utf-8
% v3.0 (Jun. 11, 2019)

\documentclass[12pt, a4paper]{ntuthesis}

% 除非校方修改了論文格式 (margins, header, footer, 浮水印, 中文數字之章別)
% 或者需要增加所用的 LaTeX 套件，
% 或者要改預設中文字型、編碼
% 否則毋須修改本檔內容

%%%%%%%%%%%%%%%%%%%%%%%%%%%%%%
\usepackage{mathtools} % 各式 AMS 數學功能
\usepackage{amssymb} % 各式 AMS 數學符號
\usepackage{mathrsfs} %草寫體數學符號，在數學模式裡用 \mathscr{E} 得草寫 E
\usepackage{bm}
\usepackage{caption}
\usepackage{subcaption}
\usepackage{tabularx}
\usepackage{url}
\usepackage[usenames,dvipsnames]{color}
%\usepackage[square, comma, numbers, sort&compress]{natbib}
\usepackage[
    colorlinks, 
    linkcolor = black, 
    citecolor = red,    % black
    urlcolor  = blue,   % black
    unicode
]{hyperref}

\usepackage{siunitx}
\usepackage{tikz}
\usetikzlibrary{arrows.meta,positioning,calc,chains,scopes}
\usetikzlibrary{shapes.geometric,angles}
\usetikzlibrary{backgrounds,fit}

\usepackage{multirow,multicol,rotating}
% 插圖套件 graphicx
\usepackage{graphicx}
% font, margin, linestretch 是設在 ntuthesis.sty 裡
\usepackage{ntuthesis}


%%%%%%%%% listings / algorithm %%%%%%%%%%
\usepackage{listings} % 程式列表套件
% hyperref跟algorithm衝突，hyperref必須放在algorithm前面
\usepackage{algorithm}
\usepackage{algorithmic}

%
% listing setting
\lstset{breaklines=true,% 過長的程式行可斷行
extendedchars=false,% 中文處理不需要 extendedchars
texcl=true,% 中文註解需要有 TeX 處理過的 comment line, 所以設成 true
comment=[l]\%\%,% 以雙「百分號」做為程式中文註解的起頭標記，配合 MATLAB
basicstyle=\small,% 小號字體, 約 10 pt 大小
commentstyle=\upshape,% 預設是斜體字，會影響註解裏的英文，改用正體
%language=Octave % 會將一些 octave 指令以粗體顯示
}
%%%%%%%%%%%%%%%%%%%%%%%%%%%%%%%%%%%%%%%%%

%%%%%%%%% fancyhdr %%%%%%%%%%
% 增強功能型頁楣 / 頁腳套件
\usepackage{fancyhdr}  % 借用此套件來擺放浮水印 
% (佔用了 central header)
% 不需要浮水印的使用者仍可利用此套件，產生所需的 header, footer
%
% 啟動 fancy header/footer 套件
\pagestyle{fancy}
\fancyhead{}  % reset left, central, right header to empty
\fancyfoot[C]{\thepage} %中間 footer 擺放頁碼
\renewcommand{\headrulewidth}{0pt} % header 的直線; 0pt 則無線

% 如果不需要任何浮水印，則請把下列介於 >>> 與 <<< 之間
% 的文字行關掉 (行首    加上百分號)
%% 浮水印 >>> 
% 請用 Adobe Acrobat 加入
% 也會需要加入 doi 的浮水印
%% this file is encoded in utf-8
% v3.0 (Jun. 11, 2019)

% 如果浮水印不是全篇需要，請把下列介於 >>> 與 <<<
% 的「全篇浮水印專用碼」關掉 (行首加百分號)
% 參考自 Keith Reckdahl 寫的 "Using Imported Graphics in LATEX2e" (epslatex.pdf) p.39
% 如果只有特定頁需要浮水印
% 則依該頁屬性使用下列之一的命令 
% 普通頁命令 \thispagestyle{WaterMarkPage}
% plain 頁命令 \thispagestyle{PlainWaterMarkPage}
% empty 頁命令 \thispagestyle{EmptyWaterMarkPage}


% 將重複使用的浮水印章
% 圖檔是 watermark.xxx
% 副檔名可以不加，可以是 latex 系統能處裡的任何格式：pdf, gif, png, jpg, eps, ...
% 某些圖檔格式在某些工作流程可能需要作前置處裡。
% 例如，pdflatex 無法直接處理 eps 檔
%  latex + dvipdfmx 無法直接處理 pdf, gif, png, jpg, 需要用 ebb 小工具程式
%  對圖檔產生 .bb 對應檔。
%
% 寬為 5.1 cm
\newsavebox{\mywatermark}
\sbox{\mywatermark}{\includegraphics[keepaspectratio, scale=0.5]{watermark}}


% 將 central header 擺放浮水印的巨集指令
\newcommand{\PlaceWaterMark}{\fancyhead[C]{\setlength{\unitlength}{1cm}%
\begin{picture}(0,0)%
\put(4.4,-4.2){\usebox{\mywatermark}}% 圖檔擺放的位置座標
\end{picture}}%
}

\fancyhead{}  % reset left, central, right header to empty
%% 如果不需整篇論文都要浮水印
%% 則下面  >>> 與 <<< 之間的程式碼請關閉
%% >>> 全篇浮水印
\PlaceWaterMark  % 每一頁都有浮水印 (除了 plain、empty 頁以外)

% 重新定義 plain 頁面
% 每張 plain 頁面 (每一章的第一頁) 也加浮水印

\fancypagestyle{plain}{%
\fancyhead{}%
\PlaceWaterMark%
\fancyfoot{}%
\fancyfoot[C]{\thepage}
\renewcommand{\headrulewidth}{0pt}%
\renewcommand{\footrulewidth}{0pt}%
}
%% <<< 全篇浮水印

%% 如果只有一、兩頁需要有浮水印
%% 可以在該頁 (有頁碼) 使用 \thispagestyle{WaterMarkPage}
%% 此命令不影響原有的 header、footer
\fancypagestyle{WaterMarkPage}{%
\PlaceWaterMark%
}

%% 如果只有一、兩頁 plain 頁需要有浮水印 (如 摘要、自傳等)
%% 可以在該頁 (有頁碼) 使用 \thispagestyle{PlainWaterMarkPage}
%% 只有頁碼與浮水印，沒有其他的 header、footer
%% 等同於 plain page style + water mark
\fancypagestyle{PlainWaterMarkPage}{%
\fancyhead{}%
\PlaceWaterMark%
\fancyfoot{}%
\fancyfoot[C]{\thepage}
\renewcommand{\headrulewidth}{0pt}%
\renewcommand{\footrulewidth}{0pt}%
}

%% 如果只有一、兩頁 empty 頁需要有浮水印 (如封面、書名頁)
%% 可以在該頁 (無頁碼) 使用 \thispagestyle{EmptyWaterMarkPage}
%% 等同於 empty page style + water mark
\fancypagestyle{EmptyWaterMarkPage}{%
\fancyhead{}%
\PlaceWaterMark%
\fancyfoot{}%
\renewcommand{\headrulewidth}{0pt}%
\renewcommand{\footrulewidth}{0pt}%
}

%% <<< 浮水印

% 如需額外的頁楣 (header) 或 footer，請在 headerfooter.tex 裡依例修改
% 它的預設內容是都關掉，可依需要打開
% this file is encoded in utf-8
% v3.0 (Jun. 11, 2019)

%%%%%%% 其他的 header (left, right) 定義
% 底下定義了一些常見的 header 型式
% 預設情況是關掉的
% 使用者可以視需要將之打開
% 也就是把下列介於 >>> 與 <<< 之間
% 的文字行打開 (行首去掉百分號)

%% header >>>
%\renewcommand{\chaptermark}[1]{%
%\markboth{\prechaptername\ \thechapter\ \postchaptername%
%\ #1}{}%
%}  %定義 header 使用的「章」層級的戳記
%\fancyhead[L]{} % 左 header 為空
%\fancyhead[R]{\leftmark}  % 右 header 擺放「章」層級的戳記 (以 \leftmark 叫出)
%\renewcommand{\headrulewidth}{0.4pt}  % header 的直線 0.4pt; 0pt 則無線
%% <<< header

%%%%%%% 其他的 footer (left, right) 定義
% 底下定義了一些常見的 footer 型式
% 預設情況是關掉的
% 使用者可以視需要將之打開
% 也就是把下列介於 >>> 與 <<< 之間
% 的文字行打開 (行首去掉百分號)

%% footer >>>
%\fancyfoot[L]{} % 左 footer 為空
%\fancyfoot[R]{\small{NTU \LaTeX\ v3.0}} % 右 footer 擺放論文格式版本
%\renewcommand{\footrulewidth}{0.4 pt} % footer 的直線 0.4pt; 0pt 則無線
%% <<< footer

%%%%%%%%%%%%%%%%%%%%%%%%%%%%%

% 常見的 notation 指令
% 請修改 notations.tex 加入自己需要的 notations
% %%%% Common Symbol %%%%
\newcommand*{\Nb}{\mathbb{N}}
\newcommand*{\Zb}{\mathbb{Z}}
\newcommand*{\Qb}{\mathbb{Q}}
\newcommand*{\Rb}{\mathbb{R}}
\newcommand*{\Cb}{\mathbb{C}}
\newcommand*{\Eb}{\mathbb{E}}

\newcommand*{\Ac}{\mathcal{A}}
\newcommand*{\Dc}{\mathcal{D}}
\newcommand*{\Oc}{\mathcal{O}}
\newcommand*{\Tc}{\mathcal{T}}
\newcommand*{\Uc}{\mathcal{U}}
\newcommand*{\Vc}{\mathcal{V}}
\newcommand*{\Xc}{\mathcal{X}}
\newcommand*{\Yc}{\mathcal{Y}}
\newcommand*{\Zc}{\mathcal{Z}}
\newcommand*{\trans}{\mathsf{T}}

%%%% mapping Symbol %%%%
\newcommand*\bij{\lhook\joinrel\twoheadrightarrow}
\newcommand*\oneto{\hookrightarrow}
\newcommand*\onto{\twoheadrightarrow}
\newcommand*\isoto{\xrightarrow{\sim}}
\newcommand*\acts{\curvearrowright}
\newcommand*\revacts{\curvearrowleft}

%%%% set definition %%%%
% just to make sure it exists
\providecommand\given{}
% can be useful to refer to this outside \Set
\newcommand*\SetSymbol[1][]{%
  \nonscript\:#1\vert
  \allowbreak
  \nonscript\:
\mathopen{}}
\DeclarePairedDelimiterX\Set[1]\{\}{%
  \renewcommand\given{\SetSymbol[\delimsize]}
  \,#1\,
}

%%%% Probability and Statistics Notations %%%%
\DeclarePairedDelimiterXPP{\KL}[2]{D_\text{KL}}\lbrack\rbrack{}{{#1} \delimsize\Vert {#2}}
\DeclarePairedDelimiterXPP{\Prob}[1]{\Pb}\lbrace\rbrace{}{#1}
\DeclarePairedDelimiterXPP{\Ev}[1]{\Eb}\lbrack\rbrack{}{#1}
\DeclarePairedDelimiterXPP{\Evr}[2]{\Eb_{#1}}\lbrack\rbrack{}{#2}

%%%%%%%%symbol and function settings%%%%%%%%%
\DeclarePairedDelimiter{\abs}{\lvert}{\rvert}
\DeclarePairedDelimiter{\norm}{\lVert}{\rVert}
\DeclarePairedDelimiter{\inpd}{\langle}{\rangle} % inner product
\DeclarePairedDelimiter{\ceil}{\lceil}{\rceil}
\DeclarePairedDelimiter{\floor}{\lfloor}{\rfloor}
\DeclareMathOperator*{\argmin}{arg\,min}
\DeclareMathOperator*{\argmax}{arg\,max}
\newcommand*{\bv}[1]{\mathbf{#1}}

%% some function name
\DeclareMathOperator{\ELBO}{ELBO}
\DeclareMathOperator{\ReLU}{ReLU}



% Title Page
% 請修改 title_settings.tex 設定論文題目/指導教授等變數
\renewcommand{\enTitle}{Research of Discrete Speech Representations}  %英文標題
\renewcommand{\zhTitle}{語音離散化表徵}  %中文標題
\renewcommand{\authorZhName}{陳建成}  %作者中文姓名
\renewcommand{\authorEnName}{Chien-Cheng Chen}  %作者英文姓名
\renewcommand{\authorStudentID}{R09942097}  %作者學號
\renewcommand{\advisorZhName}{李琳山}  %指導教授中文姓名
\renewcommand{\advisorEnName}{Lin-shan Lee}  %指導教授英文姓名
\renewcommand{\zhCollegeName}{電機資訊學院}  %學院中文名稱
\renewcommand{\enCollegeName}{College of Electrical Engineering and Computer Science}  %學院英文名稱
\renewcommand{\zhDepartmentName}{電信工程學研究所}  %系所中文名稱
\renewcommand{\enDepartmentName}{Graduate Institute of Communication Engineering}  %系所英文名稱
\renewcommand{\rocYear}{一百一十三}  %中華民國紀年年份
\renewcommand{\zhMonth}{六}  %中文月份
\renewcommand{\enYear}{2024}  %公元紀年
\renewcommand{\enMonth}{June}  %英文月份
\renewcommand{\oralDate}{113 年 6 月 1 日}  %口試日期


% 載入中文名詞的定義：例如，Figure -->「圖」, Chapter -->「第 x 章」
% this file is encoded in utf-8
% v3.0 (Jun. 11, 2019)

% 下列中文名詞的定義，如果以註解方式關閉取消，
% 則會以系統原先的預設值 (英文) 替代
% 名詞 \prechaptername 預設值為 Chapter
% 名詞 \postchaptername 預設值為空字串
% 名詞 \tablename 預設值為 Table
% 名詞 \figurename 預設值為 Figure
\renewcommand\prechaptername{第} % 出現在每一章的開頭的「第 x 章」
\renewcommand\postchaptername{章}
\renewcommand{\tablename}{表} % 在文章中 table caption 會以「表 x」表示
\renewcommand{\figurename}{圖} % 在文章中 figure caption 會以「圖 x」表示

% 下列中文名詞的定義，用於論文固定的各部分之命名 (出現於目錄與該頁標題)
\newcommand{\nameInnerCover}{書名頁}
\newcommand{\nameCommitteeForm}{論文口試委員審定書}
\newcommand{\nameCopyrightForm}{授權書}
\newcommand{\nameCabstract}{中文摘要}
\newcommand{\nameEabstract}{英文摘要}
\newcommand{\nameAckn}{誌謝}
\newcommand{\nameToc}{目錄}
\renewcommand{\contentsname}{目錄}
\newcommand{\nameLot}{表目錄}
\newcommand{\nameTof}{圖目錄}
\newcommand{\nameSlist}{符號說明}
\newcommand{\nameRef}{參考文獻}
\newcommand{\nameVita}{自傳}


% 如果不需要以中文數字一、二、三呈現章別，例如「第一章」
% 則請把下列介於 >>> 與 <<< 之間
% 的文字行關掉 (行首加上百分號), 會以「第 1 章」呈現
%% 中文數字章別 >>>
% this file is encoded in utf-8
% v3.0 (Jun. 11, 2019)

% 請依需要選擇其中一種表現方式，把它所對應的指令列打開，其他沒有用到的表現方式的對應指令列請關閉。(用行首百分號)

%% 第一種目錄格式：
%%	1  簡介 ............................ 1
%%
%%      章別 (chapter counter) 「1」前後沒有其他文字，
%%
%%      內文章標題是
%%		第 1 章	簡介
%%	\tocprechaptername, \tocpostchaptername 都設成沒有內容的空字串
%%	\tocChNumberWidth 設成 1.4em (預設)
%%      底下三行指令請打開
%\renewcommand\tocprechaptername{}
%\renewcommand\tocpostchaptername{}
%\setlength{\tocChNumberWidth}{1.4em}


%% 第二種目錄格式：
%%	一、簡介 ............................ 1
%%
%%      章別 (chapter counter) 「一」前沒有文字，後有頓號，
%%
%%      內文章標題是
%%		第一章		簡介
%%	\tocprechaptername 設成沒有內容的空字串
%%	\tocpostchaptername 設成頓號
%%	\tocChNumberWidth 設成 2em
%%      底下四行指令請打開 (預設)
\renewcommand\countermapping[1]{\zhnumber{#1}}
\renewcommand\tocprechaptername{}
\renewcommand\tocpostchaptername{、}
\setlength{\tocChNumberWidth}{2em}


%% 第三種目錄格式：
%%	第一章、簡介 ......................... 1
%%
%%      章別 (chapter counter) 「一」前有「第」，後有「章」與頓號，
%%      內文章標題是
%%		第一章		簡介
%%	\tocprechaptername 設成「第」
%%	\tocpostchaptername 設成「章、」
%%	\tocChNumberWidth 設成 3em
%%      底下四行指令請打開
%\renewcommand\countermapping[1]{\CJKnumber{#1}}
%\renewcommand\tocprechaptername{第}
%\renewcommand\tocpostchaptername{章、}
%\setlength{\tocChNumberWidth}{3em}



%% 可以依照需要作彈性的設定
%%
%% 章別 (數字，包括後面的字串) 的寬度 \tocChNumberWidth，
%% 會影響章名與章別之間的間隔 (太少則相疊，太多則留白)
%% 建議設成 \tocpostchaptername 內容字數加一，做為 em 的倍數，
%% 但至少也要有 1.4 倍。

%% <<< 中文數字章別

%%%%%%%%%%%%%%%%%%%%%%%%%%%%%
%  end of preamble
%%%%%%%%%%%%%%%%%%%%%%%%%%%%%
%
\begin{document}
\CJKindent  %%% ZZZ %%%  段首內縮兩格

%%% 以下載入前頁、本文、後頁

% \NTUtitlepage  % 產生論文封面

\newpage
\setcounter{page}{1}
\pagenumbering{roman}

\NTUoralpage  % 產生口試委員會審定書

\mydoublespacing
\begin{acknowledgement} %誌謝
    兩年來的碩士生活，隨著這份論文的誕生而告了一個段落。回思過去兩年來的生活，有做出成果的喜悅、也有處處碰壁時的苦悶、也有跟同學共同奮鬥的記憶，自己在這兩年中實在成長了許多，而這些都要感謝實驗室的大家長：李琳山教授。教授在實驗室營造了自由研究的氛圍，讓實驗室的同學都能按自己喜好自由發展研究方向，並從旁關心協助同學的研究。我在這樣的氛圍下也受益許多，學習到了許多做研究與做人處事的方法。

    這些研究能夠順利完成要由衷地感謝我的家人，他們無論何時都支持我的決定，並且從旁給我協助，在我最忙碌而都很晚回家的那段時間，他們也是很包容我，並給予我支持。

    與實驗室的同學相處的這段時光將是碩士生活中最難忘的一段日子，我們無論是修課或是研究上都是彼此的好戰友，我從你們身上都學到很多；下一屆的你們將是實驗室下一代的主力，祝你們未來研究順利！

    最後要感謝我的朋友、同學、以及伙伴們，不論是平常一起吃飯聊天、有正事時的一起奮鬥、或一起出去玩，你們都是我平常生活上最大的支持！

\end{acknowledgement}

\begin{zhAbstract} % 中文摘要
    這是中文摘要
\end{zhAbstract}

\begin{enAbstract} % 英文摘要
    This is an English abstract.
\end{enAbstract}

{
%\zhKaiFont
\mysinglespacing\selectfont
\tableofcontents %目錄

\listoffigures  %圖目錄

\listoftables  %表目錄
\par
}

\newpage
\setcounter{page}{1}
\pagenumbering{arabic}
  % omitted as...: 
% <<<
% \NTUtitlepage  % 產生論文封面

% \newpage
\setcounter{page}{1}
\pagenumbering{roman}

% \NTUoralpage  % 產生口試委員會審定書

\mydoublespacing
%% \begin{acknowledgement} %誌謝
%%     我不想看到

%% \end{acknowledgement}%% 

%% \begin{zhAbstract} % 中文摘要
%%     這是中文摘要
%% \end{zhAbstract}%% 

%% \begin{enAbstract} % 英文摘要
%%     This is an English abstract.
%% \end{enAbstract}

%%%%%% 先藏起來
%%%%%%%% {
%%%%%%%% %\zhKaiFont
%%%%%%%% \mysinglespacing\selectfont
%%%%%%%% \tableofcontents %目錄
%%%%%%%% 
%%%%%%%% % \listoffigures  %圖目錄
%%%%%%%% 
%%%%%%%% % \listoftables  %表目錄
%%%%%%%% \par
%%%%%%%% }
%%%%%% <<<

\newpage
\setcounter{page}{1}
\pagenumbering{arabic}

% >>>
% this file is encoded in utf-8
% v3.0 (Jun. 11, 2019)

\documentclass[12pt, a4paper]{ntuthesis}

% 除非校方修改了論文格式 (margins, header, footer, 浮水印, 中文數字之章別)
% 或者需要增加所用的 LaTeX 套件，
% 或者要改預設中文字型、編碼
% 否則毋須修改本檔內容

%%%%%%%%%%%%%%%%%%%%%%%%%%%%%%
\usepackage{mathtools} % 各式 AMS 數學功能
\usepackage{amssymb} % 各式 AMS 數學符號
\usepackage{mathrsfs} %草寫體數學符號，在數學模式裡用 \mathscr{E} 得草寫 E
\usepackage{bm}
\usepackage{caption}
\usepackage{subcaption}
\usepackage{tabularx}
\usepackage{url}
\usepackage[usenames,dvipsnames]{color}
%\usepackage[square, comma, numbers, sort&compress]{natbib}
\usepackage[
    colorlinks, 
    linkcolor = black, 
    citecolor = red,    % black
    urlcolor  = blue,   % black
    unicode
]{hyperref}

\usepackage{siunitx}
\usepackage{tikz}
\usetikzlibrary{arrows.meta,positioning,calc,chains,scopes}
\usetikzlibrary{shapes.geometric,angles}
\usetikzlibrary{backgrounds,fit}

\usepackage{multirow,multicol,rotating}
% 插圖套件 graphicx
\usepackage{graphicx}
% font, margin, linestretch 是設在 ntuthesis.sty 裡
\usepackage{ntuthesis}


%%%%%%%%% listings / algorithm %%%%%%%%%%
\usepackage{listings} % 程式列表套件
% hyperref跟algorithm衝突，hyperref必須放在algorithm前面
\usepackage{algorithm}
\usepackage{algorithmic}

%
% listing setting
\lstset{breaklines=true,% 過長的程式行可斷行
extendedchars=false,% 中文處理不需要 extendedchars
texcl=true,% 中文註解需要有 TeX 處理過的 comment line, 所以設成 true
comment=[l]\%\%,% 以雙「百分號」做為程式中文註解的起頭標記，配合 MATLAB
basicstyle=\small,% 小號字體, 約 10 pt 大小
commentstyle=\upshape,% 預設是斜體字，會影響註解裏的英文，改用正體
%language=Octave % 會將一些 octave 指令以粗體顯示
}
%%%%%%%%%%%%%%%%%%%%%%%%%%%%%%%%%%%%%%%%%

%%%%%%%%% fancyhdr %%%%%%%%%%
% 增強功能型頁楣 / 頁腳套件
\usepackage{fancyhdr}  % 借用此套件來擺放浮水印 
% (佔用了 central header)
% 不需要浮水印的使用者仍可利用此套件，產生所需的 header, footer
%
% 啟動 fancy header/footer 套件
\pagestyle{fancy}
\fancyhead{}  % reset left, central, right header to empty
\fancyfoot[C]{\thepage} %中間 footer 擺放頁碼
\renewcommand{\headrulewidth}{0pt} % header 的直線; 0pt 則無線

% 如果不需要任何浮水印，則請把下列介於 >>> 與 <<< 之間
% 的文字行關掉 (行首    加上百分號)
%% 浮水印 >>> 
% 請用 Adobe Acrobat 加入
% 也會需要加入 doi 的浮水印
%% this file is encoded in utf-8
% v3.0 (Jun. 11, 2019)

% 如果浮水印不是全篇需要，請把下列介於 >>> 與 <<<
% 的「全篇浮水印專用碼」關掉 (行首加百分號)
% 參考自 Keith Reckdahl 寫的 "Using Imported Graphics in LATEX2e" (epslatex.pdf) p.39
% 如果只有特定頁需要浮水印
% 則依該頁屬性使用下列之一的命令 
% 普通頁命令 \thispagestyle{WaterMarkPage}
% plain 頁命令 \thispagestyle{PlainWaterMarkPage}
% empty 頁命令 \thispagestyle{EmptyWaterMarkPage}


% 將重複使用的浮水印章
% 圖檔是 watermark.xxx
% 副檔名可以不加，可以是 latex 系統能處裡的任何格式：pdf, gif, png, jpg, eps, ...
% 某些圖檔格式在某些工作流程可能需要作前置處裡。
% 例如，pdflatex 無法直接處理 eps 檔
%  latex + dvipdfmx 無法直接處理 pdf, gif, png, jpg, 需要用 ebb 小工具程式
%  對圖檔產生 .bb 對應檔。
%
% 寬為 5.1 cm
\newsavebox{\mywatermark}
\sbox{\mywatermark}{\includegraphics[keepaspectratio, scale=0.5]{watermark}}


% 將 central header 擺放浮水印的巨集指令
\newcommand{\PlaceWaterMark}{\fancyhead[C]{\setlength{\unitlength}{1cm}%
\begin{picture}(0,0)%
\put(4.4,-4.2){\usebox{\mywatermark}}% 圖檔擺放的位置座標
\end{picture}}%
}

\fancyhead{}  % reset left, central, right header to empty
%% 如果不需整篇論文都要浮水印
%% 則下面  >>> 與 <<< 之間的程式碼請關閉
%% >>> 全篇浮水印
\PlaceWaterMark  % 每一頁都有浮水印 (除了 plain、empty 頁以外)

% 重新定義 plain 頁面
% 每張 plain 頁面 (每一章的第一頁) 也加浮水印

\fancypagestyle{plain}{%
\fancyhead{}%
\PlaceWaterMark%
\fancyfoot{}%
\fancyfoot[C]{\thepage}
\renewcommand{\headrulewidth}{0pt}%
\renewcommand{\footrulewidth}{0pt}%
}
%% <<< 全篇浮水印

%% 如果只有一、兩頁需要有浮水印
%% 可以在該頁 (有頁碼) 使用 \thispagestyle{WaterMarkPage}
%% 此命令不影響原有的 header、footer
\fancypagestyle{WaterMarkPage}{%
\PlaceWaterMark%
}

%% 如果只有一、兩頁 plain 頁需要有浮水印 (如 摘要、自傳等)
%% 可以在該頁 (有頁碼) 使用 \thispagestyle{PlainWaterMarkPage}
%% 只有頁碼與浮水印，沒有其他的 header、footer
%% 等同於 plain page style + water mark
\fancypagestyle{PlainWaterMarkPage}{%
\fancyhead{}%
\PlaceWaterMark%
\fancyfoot{}%
\fancyfoot[C]{\thepage}
\renewcommand{\headrulewidth}{0pt}%
\renewcommand{\footrulewidth}{0pt}%
}

%% 如果只有一、兩頁 empty 頁需要有浮水印 (如封面、書名頁)
%% 可以在該頁 (無頁碼) 使用 \thispagestyle{EmptyWaterMarkPage}
%% 等同於 empty page style + water mark
\fancypagestyle{EmptyWaterMarkPage}{%
\fancyhead{}%
\PlaceWaterMark%
\fancyfoot{}%
\renewcommand{\headrulewidth}{0pt}%
\renewcommand{\footrulewidth}{0pt}%
}

%% <<< 浮水印

% 如需額外的頁楣 (header) 或 footer，請在 headerfooter.tex 裡依例修改
% 它的預設內容是都關掉，可依需要打開
% this file is encoded in utf-8
% v3.0 (Jun. 11, 2019)

%%%%%%% 其他的 header (left, right) 定義
% 底下定義了一些常見的 header 型式
% 預設情況是關掉的
% 使用者可以視需要將之打開
% 也就是把下列介於 >>> 與 <<< 之間
% 的文字行打開 (行首去掉百分號)

%% header >>>
%\renewcommand{\chaptermark}[1]{%
%\markboth{\prechaptername\ \thechapter\ \postchaptername%
%\ #1}{}%
%}  %定義 header 使用的「章」層級的戳記
%\fancyhead[L]{} % 左 header 為空
%\fancyhead[R]{\leftmark}  % 右 header 擺放「章」層級的戳記 (以 \leftmark 叫出)
%\renewcommand{\headrulewidth}{0.4pt}  % header 的直線 0.4pt; 0pt 則無線
%% <<< header

%%%%%%% 其他的 footer (left, right) 定義
% 底下定義了一些常見的 footer 型式
% 預設情況是關掉的
% 使用者可以視需要將之打開
% 也就是把下列介於 >>> 與 <<< 之間
% 的文字行打開 (行首去掉百分號)

%% footer >>>
%\fancyfoot[L]{} % 左 footer 為空
%\fancyfoot[R]{\small{NTU \LaTeX\ v3.0}} % 右 footer 擺放論文格式版本
%\renewcommand{\footrulewidth}{0.4 pt} % footer 的直線 0.4pt; 0pt 則無線
%% <<< footer

%%%%%%%%%%%%%%%%%%%%%%%%%%%%%

% 常見的 notation 指令
% 請修改 notations.tex 加入自己需要的 notations
% %%%% Common Symbol %%%%
\newcommand*{\Nb}{\mathbb{N}}
\newcommand*{\Zb}{\mathbb{Z}}
\newcommand*{\Qb}{\mathbb{Q}}
\newcommand*{\Rb}{\mathbb{R}}
\newcommand*{\Cb}{\mathbb{C}}
\newcommand*{\Eb}{\mathbb{E}}

\newcommand*{\Ac}{\mathcal{A}}
\newcommand*{\Dc}{\mathcal{D}}
\newcommand*{\Oc}{\mathcal{O}}
\newcommand*{\Tc}{\mathcal{T}}
\newcommand*{\Uc}{\mathcal{U}}
\newcommand*{\Vc}{\mathcal{V}}
\newcommand*{\Xc}{\mathcal{X}}
\newcommand*{\Yc}{\mathcal{Y}}
\newcommand*{\Zc}{\mathcal{Z}}
\newcommand*{\trans}{\mathsf{T}}

%%%% mapping Symbol %%%%
\newcommand*\bij{\lhook\joinrel\twoheadrightarrow}
\newcommand*\oneto{\hookrightarrow}
\newcommand*\onto{\twoheadrightarrow}
\newcommand*\isoto{\xrightarrow{\sim}}
\newcommand*\acts{\curvearrowright}
\newcommand*\revacts{\curvearrowleft}

%%%% set definition %%%%
% just to make sure it exists
\providecommand\given{}
% can be useful to refer to this outside \Set
\newcommand*\SetSymbol[1][]{%
  \nonscript\:#1\vert
  \allowbreak
  \nonscript\:
\mathopen{}}
\DeclarePairedDelimiterX\Set[1]\{\}{%
  \renewcommand\given{\SetSymbol[\delimsize]}
  \,#1\,
}

%%%% Probability and Statistics Notations %%%%
\DeclarePairedDelimiterXPP{\KL}[2]{D_\text{KL}}\lbrack\rbrack{}{{#1} \delimsize\Vert {#2}}
\DeclarePairedDelimiterXPP{\Prob}[1]{\Pb}\lbrace\rbrace{}{#1}
\DeclarePairedDelimiterXPP{\Ev}[1]{\Eb}\lbrack\rbrack{}{#1}
\DeclarePairedDelimiterXPP{\Evr}[2]{\Eb_{#1}}\lbrack\rbrack{}{#2}

%%%%%%%%symbol and function settings%%%%%%%%%
\DeclarePairedDelimiter{\abs}{\lvert}{\rvert}
\DeclarePairedDelimiter{\norm}{\lVert}{\rVert}
\DeclarePairedDelimiter{\inpd}{\langle}{\rangle} % inner product
\DeclarePairedDelimiter{\ceil}{\lceil}{\rceil}
\DeclarePairedDelimiter{\floor}{\lfloor}{\rfloor}
\DeclareMathOperator*{\argmin}{arg\,min}
\DeclareMathOperator*{\argmax}{arg\,max}
\newcommand*{\bv}[1]{\mathbf{#1}}

%% some function name
\DeclareMathOperator{\ELBO}{ELBO}
\DeclareMathOperator{\ReLU}{ReLU}



% Title Page
% 請修改 title_settings.tex 設定論文題目/指導教授等變數
\renewcommand{\enTitle}{Research of Discrete Speech Representations}  %英文標題
\renewcommand{\zhTitle}{語音離散化表徵}  %中文標題
\renewcommand{\authorZhName}{陳建成}  %作者中文姓名
\renewcommand{\authorEnName}{Chien-Cheng Chen}  %作者英文姓名
\renewcommand{\authorStudentID}{R09942097}  %作者學號
\renewcommand{\advisorZhName}{李琳山}  %指導教授中文姓名
\renewcommand{\advisorEnName}{Lin-shan Lee}  %指導教授英文姓名
\renewcommand{\zhCollegeName}{電機資訊學院}  %學院中文名稱
\renewcommand{\enCollegeName}{College of Electrical Engineering and Computer Science}  %學院英文名稱
\renewcommand{\zhDepartmentName}{電信工程學研究所}  %系所中文名稱
\renewcommand{\enDepartmentName}{Graduate Institute of Communication Engineering}  %系所英文名稱
\renewcommand{\rocYear}{一百一十三}  %中華民國紀年年份
\renewcommand{\zhMonth}{六}  %中文月份
\renewcommand{\enYear}{2024}  %公元紀年
\renewcommand{\enMonth}{June}  %英文月份
\renewcommand{\oralDate}{113 年 6 月 1 日}  %口試日期


% 載入中文名詞的定義：例如，Figure -->「圖」, Chapter -->「第 x 章」
% this file is encoded in utf-8
% v3.0 (Jun. 11, 2019)

% 下列中文名詞的定義，如果以註解方式關閉取消，
% 則會以系統原先的預設值 (英文) 替代
% 名詞 \prechaptername 預設值為 Chapter
% 名詞 \postchaptername 預設值為空字串
% 名詞 \tablename 預設值為 Table
% 名詞 \figurename 預設值為 Figure
\renewcommand\prechaptername{第} % 出現在每一章的開頭的「第 x 章」
\renewcommand\postchaptername{章}
\renewcommand{\tablename}{表} % 在文章中 table caption 會以「表 x」表示
\renewcommand{\figurename}{圖} % 在文章中 figure caption 會以「圖 x」表示

% 下列中文名詞的定義，用於論文固定的各部分之命名 (出現於目錄與該頁標題)
\newcommand{\nameInnerCover}{書名頁}
\newcommand{\nameCommitteeForm}{論文口試委員審定書}
\newcommand{\nameCopyrightForm}{授權書}
\newcommand{\nameCabstract}{中文摘要}
\newcommand{\nameEabstract}{英文摘要}
\newcommand{\nameAckn}{誌謝}
\newcommand{\nameToc}{目錄}
\renewcommand{\contentsname}{目錄}
\newcommand{\nameLot}{表目錄}
\newcommand{\nameTof}{圖目錄}
\newcommand{\nameSlist}{符號說明}
\newcommand{\nameRef}{參考文獻}
\newcommand{\nameVita}{自傳}


% 如果不需要以中文數字一、二、三呈現章別，例如「第一章」
% 則請把下列介於 >>> 與 <<< 之間
% 的文字行關掉 (行首加上百分號), 會以「第 1 章」呈現
%% 中文數字章別 >>>
% this file is encoded in utf-8
% v3.0 (Jun. 11, 2019)

% 請依需要選擇其中一種表現方式，把它所對應的指令列打開，其他沒有用到的表現方式的對應指令列請關閉。(用行首百分號)

%% 第一種目錄格式：
%%	1  簡介 ............................ 1
%%
%%      章別 (chapter counter) 「1」前後沒有其他文字，
%%
%%      內文章標題是
%%		第 1 章	簡介
%%	\tocprechaptername, \tocpostchaptername 都設成沒有內容的空字串
%%	\tocChNumberWidth 設成 1.4em (預設)
%%      底下三行指令請打開
%\renewcommand\tocprechaptername{}
%\renewcommand\tocpostchaptername{}
%\setlength{\tocChNumberWidth}{1.4em}


%% 第二種目錄格式：
%%	一、簡介 ............................ 1
%%
%%      章別 (chapter counter) 「一」前沒有文字，後有頓號，
%%
%%      內文章標題是
%%		第一章		簡介
%%	\tocprechaptername 設成沒有內容的空字串
%%	\tocpostchaptername 設成頓號
%%	\tocChNumberWidth 設成 2em
%%      底下四行指令請打開 (預設)
\renewcommand\countermapping[1]{\zhnumber{#1}}
\renewcommand\tocprechaptername{}
\renewcommand\tocpostchaptername{、}
\setlength{\tocChNumberWidth}{2em}


%% 第三種目錄格式：
%%	第一章、簡介 ......................... 1
%%
%%      章別 (chapter counter) 「一」前有「第」，後有「章」與頓號，
%%      內文章標題是
%%		第一章		簡介
%%	\tocprechaptername 設成「第」
%%	\tocpostchaptername 設成「章、」
%%	\tocChNumberWidth 設成 3em
%%      底下四行指令請打開
%\renewcommand\countermapping[1]{\CJKnumber{#1}}
%\renewcommand\tocprechaptername{第}
%\renewcommand\tocpostchaptername{章、}
%\setlength{\tocChNumberWidth}{3em}



%% 可以依照需要作彈性的設定
%%
%% 章別 (數字，包括後面的字串) 的寬度 \tocChNumberWidth，
%% 會影響章名與章別之間的間隔 (太少則相疊，太多則留白)
%% 建議設成 \tocpostchaptername 內容字數加一，做為 em 的倍數，
%% 但至少也要有 1.4 倍。

%% <<< 中文數字章別

%%%%%%%%%%%%%%%%%%%%%%%%%%%%%
%  end of preamble
%%%%%%%%%%%%%%%%%%%%%%%%%%%%%
%
\begin{document}
\CJKindent  %%% ZZZ %%%  段首內縮兩格

%%% 以下載入前頁、本文、後頁

% \NTUtitlepage  % 產生論文封面

\newpage
\setcounter{page}{1}
\pagenumbering{roman}

\NTUoralpage  % 產生口試委員會審定書

\mydoublespacing
\begin{acknowledgement} %誌謝
    兩年來的碩士生活，隨著這份論文的誕生而告了一個段落。回思過去兩年來的生活，有做出成果的喜悅、也有處處碰壁時的苦悶、也有跟同學共同奮鬥的記憶，自己在這兩年中實在成長了許多，而這些都要感謝實驗室的大家長：李琳山教授。教授在實驗室營造了自由研究的氛圍，讓實驗室的同學都能按自己喜好自由發展研究方向，並從旁關心協助同學的研究。我在這樣的氛圍下也受益許多，學習到了許多做研究與做人處事的方法。

    這些研究能夠順利完成要由衷地感謝我的家人，他們無論何時都支持我的決定，並且從旁給我協助，在我最忙碌而都很晚回家的那段時間，他們也是很包容我，並給予我支持。

    與實驗室的同學相處的這段時光將是碩士生活中最難忘的一段日子，我們無論是修課或是研究上都是彼此的好戰友，我從你們身上都學到很多；下一屆的你們將是實驗室下一代的主力，祝你們未來研究順利！

    最後要感謝我的朋友、同學、以及伙伴們，不論是平常一起吃飯聊天、有正事時的一起奮鬥、或一起出去玩，你們都是我平常生活上最大的支持！

\end{acknowledgement}

\begin{zhAbstract} % 中文摘要
    這是中文摘要
\end{zhAbstract}

\begin{enAbstract} % 英文摘要
    This is an English abstract.
\end{enAbstract}

{
%\zhKaiFont
\mysinglespacing\selectfont
\tableofcontents %目錄

\listoffigures  %圖目錄

\listoftables  %表目錄
\par
}

\newpage
\setcounter{page}{1}
\pagenumbering{arabic}
  % omitted as...: 
% <<<
% \NTUtitlepage  % 產生論文封面

% \newpage
\setcounter{page}{1}
\pagenumbering{roman}

% \NTUoralpage  % 產生口試委員會審定書

\mydoublespacing
%% \begin{acknowledgement} %誌謝
%%     我不想看到

%% \end{acknowledgement}%% 

%% \begin{zhAbstract} % 中文摘要
%%     這是中文摘要
%% \end{zhAbstract}%% 

%% \begin{enAbstract} % 英文摘要
%%     This is an English abstract.
%% \end{enAbstract}

%%%%%% 先藏起來
%%%%%%%% {
%%%%%%%% %\zhKaiFont
%%%%%%%% \mysinglespacing\selectfont
%%%%%%%% \tableofcontents %目錄
%%%%%%%% 
%%%%%%%% % \listoffigures  %圖目錄
%%%%%%%% 
%%%%%%%% % \listoftables  %表目錄
%%%%%%%% \par
%%%%%%%% }
%%%%%% <<<

\newpage
\setcounter{page}{1}
\pagenumbering{arabic}

% >>>
% !TeX root = ./main.tex

% !TeX root = ../../../main.tex

\newcommand{\NOTEHERE}[1]{\colorbox{yellow}{\textcolor{blue}{
    \lparen\lparen\lparen}\textcolor{
        red}{#1}\textcolor{blue}{\rparen\rparen\rparen}}}

% newline
\newcommand{\PROBABLY}[1]{}
%% \newcommand{\PROBABLY}[1]{\colorbox{red}{\textcolor{green}{
%%     \lbrack\lbrack\lbrack}\textcolor{green}{#1}\textcolor{
%%         green}{\rbrack\rbrack\rbrack}}}

% \textcolor{green}{#1}
\newcommand{\FIXTHIS}[1]{\par\makebox[\textwidth]{
    \colorbox{pink}{\textcolor{green}{
        \lbrack\lbrack\lbrack}\textcolor{
            blue}{#1}\textcolor{green}{\rbrack\rbrack\rbrack}}}}

\newcommand{\mysetcounter}[2]{\setcounter{#1}{\the\numexpr#2 - 1\relax}}
\newcommand{\mychcnt}[1]{\mysetcounter{chapter}{#1}}
% \newcommand{\mychcnt}[1]{}

\newcommand{\mysccnt}[1]{\mysetcounter{section}{#1}}

% \newcommand{\CITEME}{\colorbox{red}{\textcolor{green}{
%     \lparen\lparen\lparen}\textcolor{
%         white}{CITEME\textcolor{green}{\rparen\rparen\rparen}}}\:}
\newcommand{\CITEME}{\colorbox{red}{\textcolor{green}{
    \lparen}\textcolor{
        white}{?\textcolor{green}{\rparen}}}\:}




% % !TeX root = ../thesis.tex

\input{chapters/chapter0}
\input{chapters/chapter1}
\input{chapters/chapter2}
\input{chapters/chapter3}
\input{chapters/chapter4}
\input{chapters/chapter5}

我要用這樣的 citation： \cite{mcculloch1943logical}


% % \tableofcontents
% \input{1}
% \input{2}
% \input{3}
\mychcnt{4}
\input{4}
% \input{5}


\tableofcontents

\newpage
\

\input{scripts/before/1}
\input{scripts/before/2}
\input{scripts/before/3}
\input{scripts/before/4}
% \mychcnt{5}
\input{scripts/before/5}
% \newpage
% \



% \newpage
\ % The empty page
\newpage
% \newpage
% \
% \newpage

% % this file is encoded in utf-8
% v3.0 (Jun. 11, 2019)

%%% 參考文獻
\newpage
\phantomsection % for hyperref to register this
\addcontentsline{toc}{chapter}{\nameRef}
\renewcommand{\bibname}{\protect\makebox[5cm][s]{\nameRef}}
%  \makebox{} is fragile; need protect
\bibliographystyle{IEEEtran}  % 使用 IEEE Trans 期刊格式
\bibliography{thesis}


%%% 附錄
%\input{my_appendix.tex}

%%% 自傳
%\newpage
%\chapter*{\protect\makebox[5cm][s]{\nameVita}} % \makebox{} is fragile; need protect
%\phantomsection % for hyperref to register this
%\addcontentsline{toc}{chapter}{\nameVita}
%\input{my_vita.tex}



% <<<
\newpage
\phantomsection % for hyperref to register this
\addcontentsline{toc}{chapter}{\nameRef}
\renewcommand{\bibname}{\protect\makebox[3.5cm][s]{\nameRef}}  % fixed by JEFF
% \renewcommand{\bibname}{\protect\makebox[5cm][s]{\nameRef}}  % original
%  \makebox{} is fragile; need protect
\bibliographystyle{IEEEtran}  % 使用 IEEE Trans 期刊格式
\bibliography{citations,scripts/zz}

% >>>

\end{document}

% % this file is encoded in utf-8
% v3.0 (Jun. 11, 2019)

%%% 參考文獻
\newpage
\phantomsection % for hyperref to register this
\addcontentsline{toc}{chapter}{\nameRef}
\renewcommand{\bibname}{\protect\makebox[5cm][s]{\nameRef}}
%  \makebox{} is fragile; need protect
\bibliographystyle{IEEEtran}  % 使用 IEEE Trans 期刊格式
\bibliography{thesis}


%%% 附錄
%\input{my_appendix.tex}

%%% 自傳
%\newpage
%\chapter*{\protect\makebox[5cm][s]{\nameVita}} % \makebox{} is fragile; need protect
%\phantomsection % for hyperref to register this
%\addcontentsline{toc}{chapter}{\nameVita}
%\input{my_vita.tex}



% <<<
\newpage
\phantomsection % for hyperref to register this
\addcontentsline{toc}{chapter}{\nameRef}
\renewcommand{\bibname}{\protect\makebox[3.5cm][s]{\nameRef}}  % fixed by JEFF
% \renewcommand{\bibname}{\protect\makebox[5cm][s]{\nameRef}}  % original
%  \makebox{} is fragile; need protect
\bibliographystyle{IEEEtran}  % 使用 IEEE Trans 期刊格式
\bibliography{thesis}

% >>>

\end{document}
